\scanpage{24}

lieu d'expliciter

a) la correspondance entre \inserted{les $8$} sommets du cube, et les \inserted{$8$} éléments d'ordre $3$ de $G$ (``rotations-sommets'')

b) la \deleted{\unsure{relation|éléments}} \inserted{corr.} entre les $6$ paires d'arêtes antipodiques du cube, et les $6$ \deleted{\unsure{réfle|réfl}} réflexions de $G\unsure{=|\simeq}\unsure{\mathfrak S|G}_I$

c) la corr. entre les $6$ \deleted{éléments de $\operatorname{Circ}(I)$, et} \inserted{faces du cube, et} les \deleted{$6$ \unsure{autres} éléments de $G$} \inserted{permutations circulaires} (él. d'ordre $4$ de $G\simeq\operatorname{Aut}^+(\text{Cube})$) correspondants,

et de décrire également une \uline{orientation canonique} du cube, ou ce qui revient -- $=$, \unsure{une|un} \deleted{trivialisat$^{\mathrm n}$} iso.
\[
\underset{\substack{\Vert\,\mathrm{déf}\\\underline\omega_J\,(\simeq\underline\omega_I)}}{\underline\omega_{\deleted{\unsure{I|J|Q}}}}
\simeq\bigwedge_{i\in J}
\underset{\substack{\text{ens des}\\\text{deux }\deleted{\unsure{\text{élémt}}}\ \inserted{\text{autom.}}\text{ d'ordre }4\\
\text{de la structure carrée}\\Q_i\text{ définie par }i\in J\\
\text{i.e. ens des deux orientations}\\\text{de }Q_i}}{\underbrace{\underline\omega_i}}
\]

Comme on a
\[
\underline\omega_i\simeq
\underset{\substack{\deleted{\unsure{\text{ou}}}\,\text{ens.}\\\text{des deux}\\
\text{codiagonales}\\\deleted{\unsure{\text{du carré}}}\ \inserted{\text{de }Q_i}}}{\underbrace{\underline\omega}}
\wedge
\underset{\substack{\text{ens des diag.}\\\text{pour le carré }Q_i,\ \text{i.e. }D_i=I/\underset{\substack{\text{translation}\\\text{indicée par }i}}{\underbrace{\unsure{\underline a|\underline\sigma}_i}}}}{\underbrace{D_i}}
\]
la formule s'écrit aussi
\[
\underset{\substack{\deleted{\unsure{\text{ens}}}\\\text{torseur}\\\text{sous }\mathbb F_2^J}}{\underbrace{\bigwedge_{i\in J}D_i}}\simeq\mathord{1\mkern-4mu\rule{.055em}{.66em}}
\]

Or on a \deleted{la suite exacte} l'inclusion
\diagram{\begin{tikzcd}[ampersand replacement=\&,row sep=0pt]
I\arrow[r,hook]\&\displaystyle\prod_{i\in J}D_i\\
\substack{\text{torseur}\\\text{sous }V_J(\mathbb F_2)}\&\substack{\subset\ \text{torseur sous}\\\mathbb F_2^J}
\end{tikzcd}}

% archive sheet number: 23

